\documentclass[12pt,a4paper]{article}

% --- Packages ---
\usepackage[utf8]{inputenc}
\usepackage[T1]{fontenc}
\usepackage[french]{babel}
\usepackage{amsmath, amssymb, amsfonts}
\usepackage{graphicx}
\usepackage{hyperref}
\usepackage{geometry}
\usepackage{listings}
\usepackage{xcolor}
\usepackage{titlesec}
\usepackage{fancyhdr}

% --- Configuration ---
\geometry{margin=2.5cm}
\hypersetup{
    colorlinks=true,
    linkcolor=blue,
    filecolor=magenta,      
    urlcolor=cyan,
    pdftitle={Rapport Mini Projet : IA et Cryptographie},
}

\definecolor{codegreen}{rgb}{0,0.6,0}
\definecolor{codegray}{rgb}{0.5,0.5,0.5}
\definecolor{codepurple}{rgb}{0.58,0,0.82}
\definecolor{backcolour}{rgb}{0.95,0.95,0.92}

\lstdefinestyle{mystyle}{
    backgroundcolor=\color{backcolour},   
    commentstyle=\color{codegreen},
    keywordstyle=\color{magenta},
    numberstyle=\tiny\color{codegray},
    stringstyle=\color{codepurple},
    basicstyle=\ttfamily\footnotesize,
    breakatwhitespace=false,         
    breaklines=true,                 
    captionpos=b,                    
    keepspaces=true,                 
    numbers=left,                    
    numbersep=5pt,                  
    showspaces=false,                
    showstringspaces=false,
    showtabs=false,                  
    tabsize=2
}
\lstset{style=mystyle}

% --- Page de Garde ---
\begin{document}

\begin{titlepage}
    \centering
    \includegraphics[width=0.4\textwidth]{example-image-a} % Remplacez par le logo de l'ENSA si disponible
    \vspace{1cm}
    
    {\large \textbf{Ecole Nationale des Sciences Appliquées - Béni Mellal}} \\
    {\small Filière : Cycle d'Ingénieur IACS (S2)} \\
    \vspace{2cm}
    
    {\Huge \textbf{Mini-Projet : Cryptographie}} \\
    \vspace{1cm}
    {\Large \textbf{Sujet : IA pour l'optimisation de Lattice-based Cryptography}} \\
    \vspace{2cm}
    
    \begin{minipage}{0.4\textwidth}
        \begin{flushleft} \large
            \emph{Réalisé par :}\\
            \textbf{Haytham KENNOUZ} \\
            \textbf{Abdelhak ELAZZOUZI}
        \end{flushleft}
    \end{minipage}
    \hfill
    \begin{minipage}{0.4\textwidth}
        \begin{flushright} \large
            \emph{Encadré par :}\\
            \textbf{Prof. OUNACHAD}
        \end{flushright}
    \end{minipage}
    
    \vfill
    {\large Année Universitaire : 2025/2026}
\end{titlepage}

\tableofcontents
\newpage

% --- Sections ---

\section{Introduction}
La cryptographie basée sur les réseaux (\textit{Lattice-based cryptography}) représente l'une des directions les plus prometteuses pour la cryptographie post-quantique. Contrairement aux systèmes RSA ou ECC, sa sécurité repose sur des problèmes mathématiques supposés résistants même aux ordinateurs quantiques, tels que le \textit{Shortest Vector Problem} (SVP) et le \textit{Learning With Errors} (LWE).

Cependant, l'efficacité de ces systèmes et leur vulnérabilité face aux attaques dépendent fortement des algorithmes de réduction de réseaux, notamment \textbf{LLL} (\textit{Lenstra-Lenstra-Lovász}) et \textbf{BKZ} (\textit{Block-Korkine-Zolotarev}). Ce projet explore comment l'Intelligence Artificielle peut être utilisée pour optimiser ces processus complexes.

\section{Objectifs du Projet}
L'objectif principal est de concevoir un système capable de prédire les performances des algorithmes de réduction afin d'automatiser le choix des paramètres optimaux.
\begin{itemize}
    \item \textbf{Comprendre les menaces :} Analyser comment la réduction de réseau impacte la sécurité des cryptosystèmes.
    \item \textbf{Analyse Prédictive :} Utiliser le Machine Learning pour estimer le temps de calcul et la qualité de la réduction (norme du vecteur le plus court).
    \item \textbf{Optimisation :} Développer un outil capable de recommander la meilleure stratégie (LLL vs BKZ) en fonction des contraintes de temps de l'utilisateur.
\end{itemize}

\section{Méthodologie et Implémentation}

\subsection{Collecte des Données}
Nous avons utilisé la bibliothèque spécialisée \texttt{fpylll} pour générer un dataset de réseaux aléatoires de dimensions variées ($n=30$ à $60$). Pour chaque réseau, nous avons exécuté plusieurs itérations de LLL (avec différents $\delta$) et BKZ (avec différents \textit{block sizes} $\beta$).
Les données collectées incluent :
\begin{itemize}
    \item Caractéristiques : Dimension, taille des coefficients (bits), algorithme utilisé, paramètre de contrôle.
    \item Cibles : Temps d'exécution réel et norme du vecteur le plus court trouvé.
\end{itemize}

\subsection{Modèle d'Intelligence Artificielle}
Nous avons opté pour un modèle de \textbf{Random Forest Regressor} via \texttt{scikit-learn}. Ce choix se justifie par la capacité de cet algorithme à capturer les relations non-linéaires complexes entre la dimension du réseau et la complexité exponentielle des algorithmes de réduction.

\section{Résultats et Analyse}

\subsection{Performance des Modèles}
Les modèles entraînés ont montré des résultats remarquables :
\begin{itemize}
    \item \textbf{Prédiction de la Qualité :} Un coefficient de détermination $R^2 \approx 0.99$, indiquant une prédiction quasi parfaite de la norme finale du vecteur.
    \item \textbf{Prédiction du Temps :} $R^2 \approx 0.77$, ce qui est satisfaisant compte tenu de la variabilité des performances système lors de l'exécution.
\end{itemize}

\subsection{Application ALOT}
L'outil développé, \textbf{AI-driven Lattice Optimization Tool (ALOT)}, permet d'automatiser le choix stratégique. Par exemple, si un utilisateur souhaite réduire un réseau de dimension 50 en moins de 0.1 seconde, l'IA analysera toutes les combinaisons possibles et recommandera, par exemple, une réduction LLL avec $\delta=0.85$ plutôt qu'un BKZ trop coûteux.

\section{Conclusion}
Ce projet démontre que l'Intelligence Artificielle n'est pas seulement une menace pour la cryptographie (via les attaques), mais aussi un outil puissant pour son optimisation. En prédisant les performances des algorithmes de réduction, nous pouvons mieux comprendre la sécurité des réseaux et optimiser les ressources nécessaires à leur cryptanalyse ou à leur mise en œuvre sécurisée.

\section{Bibliographie}
\begin{enumerate}
    \item Lenstra, A. K., Lenstra, H. W., \& Lovász, L. (1982). \textit{Factoring polynomials with rational coefficients}.
    \item Schnorr, C. P., \& Euchner, M. (1994). \textit{Lattice basis reduction: Improved practical algorithms and solving shortest lattice vector problems}.
    \item The fpylll Team. \textit{fpylll: A Python interface for fplll}. \url{https://github.com/fplll/fpylll}.
\end{enumerate}

\end{document}
